\documentclass[11pt,a4paper]{article}
\usepackage[utf8]{inputenc}
\usepackage[italian]{babel}
\usepackage[T1]{fontenc}
\usepackage{cmap}
\usepackage[margin=2cm]{geometry}
\usepackage[table]{xcolor}
\usepackage{enumitem}
\usepackage{hyperref}
\usepackage{tcolorbox}
\usepackage{titlesec}
\usepackage{booktabs}
\usepackage{graphicx}
\usepackage{float}
\usepackage{caption}

\definecolor{primaryColor}{HTML}{0054A6}
\definecolor{secondaryColor}{HTML}{4A5568}
\definecolor{backgroundColor}{HTML}{F7FAFC}
\definecolor{borderColor}{HTML}{E2E8F0}

\hypersetup{
    colorlinks=true,
    linkcolor=primaryColor,
    urlcolor=primaryColor,
}

\titleformat{\section}
  {\color{primaryColor}\normalfont\Large\bfseries}
  {}{0em}{}
  [\color{borderColor}\titlerule]

\titleformat{\subsection}
  {\color{primaryColor}\normalfont\large\bfseries}
  {}{0em}{}

\graphicspath{{./}}

\begin{document}

\begin{tcolorbox}[colback=primaryColor,colframe=primaryColor,arc=3mm,halign=center,valign=center,height=2.5cm]
    \color{white}\LARGE\bfseries BLUEPRINT \\[0.2cm]
    \color{white!80}\large UX Design Project: Help People Help People --- TPER S.p.A. \\[0.1cm]
    \color{white!60}\normalsize Phase C --- Design Proposal
\end{tcolorbox}

\vspace{0.3cm}

\section{Introduzione}

I blueprint presentati in questa sezione definiscono i componenti per l'organizzazione dei contenuti e le connessioni tra di essi, seguendo le linee guida del corso. Sono state create tre prospettive distinte per rappresentare diversi aspetti dell'ecosistema TPER:

\begin{itemize}
    \item \textbf{BP1} --- Architettura dell'informazione di \texttt{tper.it} (prospettiva struttura)
    \item \textbf{BP2} --- Ecosistema multi-canale (prospettiva sistema)
    \item \textbf{BP3} --- User journey dal bisogno all'acquisto (prospettiva utente)
\end{itemize}

\noindent I blueprint sono stati generati con HTML+CSS per garantire una resa grafica professionale, quindi convertiti in PNG per l'inclusione in questo documento.

\vspace{0.3cm}

% =====================================================================
% BP1 — INFORMATION ARCHITECTURE
% =====================================================================
\section{BP1 --- Architettura dell'Informazione: tper.it}

\textit{Prospettiva struttura.} Questo blueprint mostra l'organizzazione gerarchica dei contenuti su \texttt{tper.it}, evidenziando la separazione netta tra contenuti informativi (gestiti direttamente sul sito) e le destinazioni di acquisto (canali esterni). Le frecce arancioni indicano i reindirizzamenti verso i canali di acquisto corretti.

\begin{figure}[H]
\centering
\includegraphics[width=\textwidth]{html/bp1-ia.png}
\caption{BP1 --- Architettura dell'Informazione di \texttt{tper.it}}
\label{fig:bp1}
\end{figure}

\subsubsection*{Legenda}
\begin{itemize}
    \item \textbf{Blu} --- Contenuto informativo gestito su \texttt{tper.it}
    \item \textbf{Arancione} --- Reindirizzamento a canale esterno per acquisto
    \item \textbf{Viola/Verde/Altri} --- Canali di acquisto esterni all'ecosistema \texttt{tper.it}
\end{itemize}

\subsubsection*{Elementi chiave}
\begin{enumerate}
    \item \textbf{Route Planner come hero}: La funzione primaria di \texttt{tper.it} è la pianificazione del viaggio. Il motore di ricerca Da--A è posizionato in evidenza nella homepage.
    \item \textbf{Niente acquisti su tper.it}: Ogni sezione che implica una transazione economica reindirizza l'utente al canale competente (solweb, Roger, Punti Tper, tabaccherie, contactless).
    \item \textbf{Separazione netta}: I contenuti informativi (tipologie, prezzi, guide) sono gestiti su \texttt{tper.it}. Le transazioni avvengono esclusivamente sui canali esterni.
    \item \textbf{Guida alla scelta}: La sezione "Dove Comprare" funge da hub decisionale che indirizza l'utente al canale più adatto alle sue esigenze.
\end{enumerate}

\vspace{0.3cm}

% =====================================================================
% BP2 — MULTI-CHANNEL ECOSYSTEM
% =====================================================================
\section{BP2 --- Ecosistema Multi-Canale TPER}

\textit{Prospettiva sistema.} Questo blueprint illustra le relazioni tra \texttt{tper.it} (canale informativo) e i 5 canali di acquisto, mostrando i meccanismi di ponte che collegano l'informazione alla transazione.

\begin{figure}[H]
\centering
\includegraphics[width=\textwidth]{html/bp2-ecosystem.png}
\caption{BP2 --- Ecosistema Multi-Canale TPER}
\label{fig:bp2}
\end{figure}

\newpage

\subsubsection*{I 5 canali di acquisto}

\begin{table}[H]
\centering
\small
\renewcommand{\arraystretch}{1.3}
\begin{tabular}{|p{2.5cm}|p{3.5cm}|p{3.5cm}|p{3.5cm}|}
\hline
\textbf{Canale} & \textbf{Tipo} & \textbf{Funzioni} & \textbf{Accesso} \\
\hline
\texttt{solweb.tper.it} & Digitale & Abbonamenti online, rinnovo, ricarica, agevolazioni ISEE & Browser web, richiede account \\
\hline
Roger App & Digitale & Biglietti digitali, acquisto, validazione, saldo contatti & Smartphone (iOS/Android) \\
\hline
Punti Tper & Fisico & Tutti i servizi, assistenza, ricariche, abbonamenti & Sportello fisico \\
\hline
Tabaccherie / PuntoLis & Fisico & Ricariche Boost, biglietti singoli & Rivendita autorizzata \\
\hline
Contactless a bordo & Ibrido & Pay-per-ride con carta o smartphone & Su veicolo TPER \\
\hline
\end{tabular}
\caption{Panoramica dei 5 canali di acquisto dell'ecosistema TPER}
\label{tab:canali}
\end{table}

\subsubsection*{Meccanismi di ponte}
\begin{itemize}
    \item \textbf{Reindirizzi web}: Link esterni da \texttt{tper.it} a \texttt{solweb.tper.it} e pagine di download Roger App
    \item \textbf{QR Code}: Scontrino rilasciato ai Punti Tper con QR code per riprendere la pratica da casa
    \item \textbf{Guide stampate}: PDF generati automaticamente con il percorso completo (R6)
    \item \textbf{Trainer-Mode}: Operatori sportello che guidano l'utente nella scelta del canale
\end{itemize}

\vspace{0.3cm}

% =====================================================================
% BP3 — USER JOURNEY
% =====================================================================
\section{BP3 --- User Journey: da tper.it all'Acquisto}

\textit{Prospettiva utente.} Questo blueprint traccia il percorso tipico di un utente HPHP dal momento in cui arriva su \texttt{tper.it} fino al completamento dell'acquisto sul canale corretto.

\begin{figure}[H]
\centering
\includegraphics[width=0.9\textwidth]{html/bp3-journey.png}
\caption{BP3 --- User Journey da \texttt{tper.it} all'Acquisto}
\label{fig:bp3}
\end{figure}

\subsubsection*{Punti di decisione}
\begin{enumerate}
    \item \textbf{Devi acquistare?}: Biforcazione tra utenti che cercano solo informazioni (orari, percorsi) e utenti che devono comprare.
    \item \textbf{Digitale o Fisico?}: Scelta basata sulle preferenze e competenze digitali dell'utente.
    \item \textbf{Hai un account?}: Per il canale digitale, distinzione tra chi ha gi\`a un account solweb e chi deve usare Roger App.
    \item \textbf{Cosa devi fare?}: Per il canale fisico, distinzione tra acquisto completo (Punti Tper), ricarica (tabaccherie) o pagamento a bordo (contactless).
\end{enumerate}

\subsubsection*{Intervento Trainer-Mode}
Il sistema attiva la modalit\`a assistita (Trainer-Mode) in due casi:
\begin{itemize}
    \item \textbf{Su richiesta}: L'utente clicca ``Aiuto'' durante il percorso
    \item \textbf{Dopo 2 errori consecutivi}: L'utente mostra difficolt\`a nella navigazione o nella scelta del canale
\end{itemize}
La Trainer-Mode fornisce: checklist dei documenti, guida alla scelta del canale, e tutorial passo-passo.

\vspace{0.5cm}

% =====================================================================
% RIEPILOGO
% =====================================================================
\section{Riepilogo dei Blueprint}

\begin{table}[H]
\centering
\small
\renewcommand{\arraystretch}{1.3}
\begin{tabular}{|p{1.8cm}|p{3cm}|p{3.5cm}|p{5cm}|}
\hline
\textbf{Blueprint} & \textbf{Prospettiva} & \textbf{Destinatari} & \textbf{Elementi coperti} \\
\hline
BP1 & Struttura & Sviluppatori, content manager & Gerarchia pagine, sezioni informative, reindirizzamenti \\
\hline
BP2 & Sistema & IT, product owner, marketing & Relazioni tra canali, meccanismi di ponte, tipologie \\
\hline
BP3 & Utente & Designer, operatori sportello & Percorso decisionale, punti di intervento, trainer-mode \\
\hline
\end{tabular}
\caption{Panoramica dei tre blueprint e loro destinatari}
\label{tab:riepilogo-bp}
\end{table}

\vspace{0.5cm}

\begin{center}
{\color{secondaryColor}\small
Documento compilato sulla base delle evidenze raccolte durante la Phase A e Phase B \\[4pt]
\textbf{Phase C --- Design Proposal: Blueprint} \\[4pt]
Progetto: UX Design Project: Help People Help People --- TPER S.p.A.}
\end{center}

\end{document}
